\documentclass{article}
\usepackage{amsmath,amssymb}
\usepackage{xurl}
\usepackage[hidelinks]{hyperref}
% Shared commands for the notes in docs/paper-gaps/.
% Every note loads this file via % Shared commands for the notes in docs/paper-gaps/.
% Every note loads this file via % Shared commands for the notes in docs/paper-gaps/.
% Every note loads this file via \input{command}. Project-wide notation
% belongs here, never in per-note redefinitions.

% --- Project configuration ------------------------------------------------
\newcommand{\projectIssueBase}{https://github.com/LionSR/QICLean/issues/}
\newcommand{\projectPrBase}{https://github.com/LionSR/QICLean/pull/}
\newcommand{\projectDocBase}{https://sirui-lu.com/QICLean/docs/}
% ---------------------------------------------------------------------------

\DeclareUnicodeCharacter{2080}{\ifmmode{}_{0}\else\textsubscript{0}\fi}
\DeclareUnicodeCharacter{2081}{\ifmmode{}_{1}\else\textsubscript{1}\fi}
\DeclareUnicodeCharacter{2082}{\ifmmode{}_{2}\else\textsubscript{2}\fi}
\DeclareUnicodeCharacter{2099}{\ifmmode{}_{n}\else\textsubscript{n}\fi}

\newcommand{\Lean}{\textsc{Lean}}

% A formal declaration name, upright sans, breakable at underscores.
\ExplSyntaxOn
\NewDocumentCommand{\leanid}{m}{
  \ifmmode
    \textsf{\detokenize{#1}}
  \else
    \group_begin:
    \urlstyle{sf}
    \tl_set:Nn \l_tmpa_tl {#1}
    \tl_replace_all:Nnn \l_tmpa_tl {\_} {_}
    \exp_args:NV \path \l_tmpa_tl
    \group_end:
  \fi
}
\ExplSyntaxOff

% Traceability links; use in footnotes, not in the running prose.
\newcommand{\ghissue}[1]{\href{\projectIssueBase#1}{GitHub issue \##1}}
\newcommand{\ghpr}[1]{\href{\projectPrBase#1}{PR \##1}}
\newcommand{\doclink}[1]{\href{\projectDocBase#1.html}{\path{#1}}}

% Domain notation, using \providecommand so a note-local refinement stays
% possible.
\providecommand{\tr}{\operatorname{tr}}

% Dirac notation and the identity operator, as in blueprint/src/macros.
\usepackage{dsfont}
\providecommand{\ket}[1]{|#1\rangle}
\providecommand{\bra}[1]{\langle #1|}
\providecommand{\braket}[2]{\langle #1 | #2 \rangle}
\providecommand{\ketbra}[2]{|#1\rangle\!\langle #2|}
\providecommand{\Id}{\mathds{1}}
\providecommand{\one}{\mathds{1}}
\providecommand{\C}{\mathbb{C}}
\providecommand{\MN}[1]{M_{#1}(\C)}
\providecommand{\MD}{M_D(\C)}

% Verdict marker: \gapnote{<kind>}{<status>}. Kind is the policy
% classification (clarification, local-correction, scope-restriction,
% unfaithful, false-source, open-gap); status is open, wip (elimination
% underway), resolved, or historical. Machine-read by the site index;
% prints a verdict line.
\newcommand{\gapnote}[2]{\par\noindent\textbf{Verdict:} #1 (#2).\par}
. Project-wide notation
% belongs here, never in per-note redefinitions.

% --- Project configuration ------------------------------------------------
\newcommand{\projectIssueBase}{https://github.com/LionSR/QICLean/issues/}
\newcommand{\projectPrBase}{https://github.com/LionSR/QICLean/pull/}
\newcommand{\projectDocBase}{https://sirui-lu.com/QICLean/docs/}
% ---------------------------------------------------------------------------

\DeclareUnicodeCharacter{2080}{\ifmmode{}_{0}\else\textsubscript{0}\fi}
\DeclareUnicodeCharacter{2081}{\ifmmode{}_{1}\else\textsubscript{1}\fi}
\DeclareUnicodeCharacter{2082}{\ifmmode{}_{2}\else\textsubscript{2}\fi}
\DeclareUnicodeCharacter{2099}{\ifmmode{}_{n}\else\textsubscript{n}\fi}

\newcommand{\Lean}{\textsc{Lean}}

% A formal declaration name, upright sans, breakable at underscores.
\ExplSyntaxOn
\NewDocumentCommand{\leanid}{m}{
  \ifmmode
    \textsf{\detokenize{#1}}
  \else
    \group_begin:
    \urlstyle{sf}
    \tl_set:Nn \l_tmpa_tl {#1}
    \tl_replace_all:Nnn \l_tmpa_tl {\_} {_}
    \exp_args:NV \path \l_tmpa_tl
    \group_end:
  \fi
}
\ExplSyntaxOff

% Traceability links; use in footnotes, not in the running prose.
\newcommand{\ghissue}[1]{\href{\projectIssueBase#1}{GitHub issue \##1}}
\newcommand{\ghpr}[1]{\href{\projectPrBase#1}{PR \##1}}
\newcommand{\doclink}[1]{\href{\projectDocBase#1.html}{\path{#1}}}

% Domain notation, using \providecommand so a note-local refinement stays
% possible.
\providecommand{\tr}{\operatorname{tr}}

% Dirac notation and the identity operator, as in blueprint/src/macros.
\usepackage{dsfont}
\providecommand{\ket}[1]{|#1\rangle}
\providecommand{\bra}[1]{\langle #1|}
\providecommand{\braket}[2]{\langle #1 | #2 \rangle}
\providecommand{\ketbra}[2]{|#1\rangle\!\langle #2|}
\providecommand{\Id}{\mathds{1}}
\providecommand{\one}{\mathds{1}}
\providecommand{\C}{\mathbb{C}}
\providecommand{\MN}[1]{M_{#1}(\C)}
\providecommand{\MD}{M_D(\C)}

% Verdict marker: \gapnote{<kind>}{<status>}. Kind is the policy
% classification (clarification, local-correction, scope-restriction,
% unfaithful, false-source, open-gap); status is open, wip (elimination
% underway), resolved, or historical. Machine-read by the site index;
% prints a verdict line.
\newcommand{\gapnote}[2]{\par\noindent\textbf{Verdict:} #1 (#2).\par}
. Project-wide notation
% belongs here, never in per-note redefinitions.

% --- Project configuration ------------------------------------------------
\newcommand{\projectIssueBase}{https://github.com/LionSR/QICLean/issues/}
\newcommand{\projectPrBase}{https://github.com/LionSR/QICLean/pull/}
\newcommand{\projectDocBase}{https://sirui-lu.com/QICLean/docs/}
% ---------------------------------------------------------------------------

\DeclareUnicodeCharacter{2080}{\ifmmode{}_{0}\else\textsubscript{0}\fi}
\DeclareUnicodeCharacter{2081}{\ifmmode{}_{1}\else\textsubscript{1}\fi}
\DeclareUnicodeCharacter{2082}{\ifmmode{}_{2}\else\textsubscript{2}\fi}
\DeclareUnicodeCharacter{2099}{\ifmmode{}_{n}\else\textsubscript{n}\fi}

\newcommand{\Lean}{\textsc{Lean}}

% A formal declaration name, upright sans, breakable at underscores.
\ExplSyntaxOn
\NewDocumentCommand{\leanid}{m}{
  \ifmmode
    \textsf{\detokenize{#1}}
  \else
    \group_begin:
    \urlstyle{sf}
    \tl_set:Nn \l_tmpa_tl {#1}
    \tl_replace_all:Nnn \l_tmpa_tl {\_} {_}
    \exp_args:NV \path \l_tmpa_tl
    \group_end:
  \fi
}
\ExplSyntaxOff

% Traceability links; use in footnotes, not in the running prose.
\newcommand{\ghissue}[1]{\href{\projectIssueBase#1}{GitHub issue \##1}}
\newcommand{\ghpr}[1]{\href{\projectPrBase#1}{PR \##1}}
\newcommand{\doclink}[1]{\href{\projectDocBase#1.html}{\path{#1}}}

% Domain notation, using \providecommand so a note-local refinement stays
% possible.
\providecommand{\tr}{\operatorname{tr}}

% Dirac notation and the identity operator, as in blueprint/src/macros.
\usepackage{dsfont}
\providecommand{\ket}[1]{|#1\rangle}
\providecommand{\bra}[1]{\langle #1|}
\providecommand{\braket}[2]{\langle #1 | #2 \rangle}
\providecommand{\ketbra}[2]{|#1\rangle\!\langle #2|}
\providecommand{\Id}{\mathds{1}}
\providecommand{\one}{\mathds{1}}
\providecommand{\C}{\mathbb{C}}
\providecommand{\MN}[1]{M_{#1}(\C)}
\providecommand{\MD}{M_D(\C)}

% Verdict marker: \gapnote{<kind>}{<status>}. Kind is the policy
% classification (clarification, local-correction, scope-restriction,
% unfaithful, false-source, open-gap); status is open, wip (elimination
% underway), resolved, or historical. Machine-read by the site index;
% prints a verdict line.
\newcommand{\gapnote}[2]{\par\noindent\textbf{Verdict:} #1 (#2).\par}

\title{The Input Scope of Semiclassical Phase Estimation}
\author{}
\date{2026-11-20}
\begin{document}
\maketitle
\gapnote{scope-restriction}{open}

\section{The source theorem}
Griffiths and Niu~\cite{GriffithsNiu1996} replace the inverse quantum Fourier
transform on a register of $t$ qubits, followed by a measurement of the whole
register in the computational basis, by measurements of the qubits one at a
time. The qubit $j$ is measured after the phase gate
$\operatorname{diag}(1,e^{-2\pi i\theta_j})$ and the Hadamard gate, where
\begin{align}
    \theta_j & = \sum_{i<j}\frac{b_i}{2^{j-i+1}}
    \label{eq:gn96_semiclassical_general_input_scope_correction}
\end{align}
is computed from the outcomes $b_i$ of the qubits measured before it. Their
claim is that the two procedures give the same distribution of outcomes for
every state of the register, and hence, inside phase estimation
(Nielsen and Chuang~\cite{NielsenChuang2010}, \S5.2), for every input state
of the control register and the target. The argument is local: in the coherent
circuit, the qubit $j$ is the control of the controlled phase gates acting on
the qubits measured after it, and these gates are diagonal in the control.
Measuring a control in the computational basis commutes with a gate that is
diagonal in it, so the measurement can be moved before the gates, which then
become phase gates on the later qubits chosen by the outcome. No property of
the input is used.

\section{The formalized statement}
The formal theorem\footnote{Formal declarations:
    \leanid{QuantumCircuit.PhaseEstimation.semiclassical_outcomeProb_eq} and
    \leanid{QuantumCircuit.PhaseEstimation.semiclassical_accurate_prob_ge} in
    \doclink{QICLean/Circuit/Algorithm/SemiclassicalPhaseEstimation}.}
concerns the whole semiclassical phase-estimation circuit, the Hadamard
gates and controlled calls of the oracle $U$ followed by the corrected
measurements. Its input is $\lvert0\cdots0\rangle\otimes u$ with
\begin{align}
    Uu & =e^{2\pi i\varphi}u,\qquad \lVert u\rVert=1,
    \label{eq:gn96_semiclassical_general_input_scope_eigenvector}
\end{align}
and its conclusion is that the readout $k$ has probability
$\lvert\alpha_k\rvert^2$, with
$\alpha_k=2^{-t}\sum_{l<2^t}e^{2\pi il(\varphi-k/2^t)}$, which is the outcome
distribution of the coherent circuit on the same input. The eigenvector
condition~\eqref{eq:gn96_semiclassical_general_input_scope_eigenvector} is a
hypothesis absent from the source's statement about the semiclassical Fourier
transform. It is the hypothesis of the phase-estimation theorem that the
semiclassical theorem is applied to, so the success probability and the number
of oracle calls of phase estimation are formalized with their source
hypotheses.

The proof uses the eigenvector input. After the calls the control register is
the product state
$\bigotimes_j2^{-1/2}(\lvert0\rangle+e^{2\pi i\varphi2^{t-1-j}}\lvert1\rangle)$
and the target is $u$. Along a branch of outcomes $b_0,\dots,b_{j-1}$ the state
stays a product, and the corrected measurement of the qubit $j$ multiplies its
factor by an explicit scalar $\beta_j$. The probability of a branch is
therefore $\prod_j\lvert\beta_j\rvert^2$, and this product is computed to be
$\lvert\alpha_k\rvert^2$. For a general input the control register is
entangled with the target after the calls, and no branch is a product state.

\section{Elimination}
The source-faithful statement is an identity of instruments that does not
mention the input: for every operator $\rho$ on the control register and the
target, the semiclassical circuit and the coherent circuit (the inverse
Fourier transform followed by the measurement of the register) have the same
outcome probabilities on $\rho$. It follows from two facts.
\begin{enumerate}
    \item \emph{Deferred measurement.} Let $P_0,P_1$ be the projections of a
          control qubit on $\lvert0\rangle$ and $\lvert1\rangle$, and let
          $W=P_0\otimes1+P_1\otimes V$ be a gate diagonal in the control. Then
          for every $\rho$ and every outcome $b$,
          \begin{align}
              (P_b\otimes1)\,W\rho W^*\,(P_b\otimes1)
               & =(1\otimes V^b)(P_b\otimes1)\rho(P_b\otimes1)(1\otimes V^b)^*.
              \label{eq:gn96_semiclassical_general_input_scope_deferred}
          \end{align}
          Hence measuring the control after $W$ is the instrument
          ``measure the control, then apply $V^b$''.
    \item \emph{The Fourier circuit.} The inverse Fourier transform is the
          product, over $j$, of the Hadamard gate on the qubit $j$ preceded by
          the controlled phase gates whose controls are the qubits $i<j$ and
          whose phases sum to $e^{-2\pi i\theta_j}$ on the outcomes $b_i$.
          Applying~\eqref{eq:gn96_semiclassical_general_input_scope_deferred}
          to every controlled phase gate, in the order of the measurements,
          turns the measurement of the transformed register into the
          semiclassical measurements, branch by branch.
\end{enumerate}
Both facts are identities of operators, so the induction runs over the
measured qubits on an arbitrary operator, with no product structure. The
eigenvector theorem is then the specialization of the general identity to the
input $\lvert0\cdots0\rangle\otimes u$, composed with the coherent
phase-estimation theorem.\footnote{Formal declaration:
    \leanid{QuantumCircuit.PhaseEstimation.outcomeProb_eq} in
    \doclink{QICLean/Circuit/Algorithm/PhaseEstimation}.}

\bibliographystyle{alpha}
\bibliography{references}
\end{document}
